\chapter{Reactiemechanismen van complexen}\label{ch:b3:complex-mechanisms}

Los een chroom(III)zout op in water dat met zuurstof-18 gemerkt is, en er
gaan dagen voorbij voordat het gemerkte water de watermoleculen heeft
vervangen die aan het metaal gebonden zijn; met een koper(II)zout is de
uitwisseling voltooid in de tijd die het mengen kost. De reactie is
dezelfde, een watermolecuul dat een metaalion verlaat en een ander dat
zijn plaats inneemt, en toch varieert haar snelheid enorm van het ene
ion tot het andere, en het aantal $d$-elektronen voorspelt een groot deel van die variatie. De manier
waarop liganden vervangen worden, bepaalt hoe platina-antikankermiddelen
gemaakt worden en hoe ze werken; de manier waarop elektronen tussen
metaalionen bewegen, bepaalt de snelheid van de ademhaling en van elke
batterij. Dit hoofdstuk behandelt de twee families van reacties van
complexen: substitutie en elektronenoverdracht, de laatste met de theorie
van Rudolph Marcus.

\begin{recall}
Het deel van jaar~2 beschreef complexen, de tandigheid van liganden,
brugliganden, isomeren $fac$ en $mer$, liganduitwisseling als elementaire stap, en
kristalveldstabilisatie-energieën met hoogspin- en
laagspinconfiguraties; het deel van jaar~1 reactiemechanismen met hun
benaderingen van de stationaire toestand en de snelheidsbepalende stap.
\Cref{ch:b3:rate-theories} gaf de \omterm{thm:b3:rate-theories:eyring}{vergelijking van Eyring} en de
betekenis van $\Delta^{\ddagger}S^\circ$, \cref{ch:b3:electrode-kinetics} de elektronenoverdracht aan een
elektrode, en \cref{ch:b3:complex-spectra-magnetism} de \omterm{def:b3:complex-spectra-magnetism:ligand-field-term}{ligandveldtermen}
en het jahn-tellereffect.
\end{recall}

\section{Labiliteit en inertie}

\begin{definition}[Labiele en inerte complexen]\label{def:b3:complex-mechanisms:labile}
Een complex is \emph{labiel}\index{labiel complex} als zijn liganden snel
vervangen worden, en een
\emph{kinetisch inert complex}\index{kinetisch inert complex}
als ze traag vervangen worden; naar Taube noemt men een complex labiel
als zijn reacties met een ligand bij \qty{0.1}{mol/L} en \qty{25}{\celsius} binnen ongeveer een
minuut afgelopen zijn. Labiliteit is een kinetische eigenschap, los van
de thermodynamische stabiliteit van het complex.
\end{definition}

De wateruitwisseling, $\mathrm{[M(H_2O)_6]^{n+} + H_2O^{\ast} \to [M(H_2O)_5(H_2O^{\ast})]^{n+} + H_2O}$, is de eenvoudigste
substitutie en legt een schaal vast. Ionen van de hoofdgroepen wisselen
sneller uit naarmate ze groter en minder geladen zijn. Onder de ionen van
de overgangsmetalen zijn die met de $d^3$-configuraties
(\ce{Cr^{3+}}) of laagspin-$d^6$ (\ce{Co^{3+}}, \ce{Rh^{3+}}, \ce{Ir^{3+}}) inert, net als de meeste
metalen van de tweede en derde rij, terwijl $d^9$ (\ce{Cu^{2+}}) en hoogspin-$d^4$ (\ce{Cr^{2+}}),
vervormd door het jahn-tellereffect, uiterst \omterm{def:b3:complex-mechanisms:labile}{labiel} zijn.

\begin{proposition}[Ligandveldactiveringsenergie]\label{prop:b3:complex-mechanisms:lfae}
Als een substitutie verloopt via een vijfgecoördineerde
vierkant-piramidale overgangstoestand, is het verlies aan
ligandveldstabilisatie op weg daarheen (de
ligandveldactiveringsenergie), in het eenvoudigste model met alleen $\sigma$-binding, het grootst voor laagspin-$d^6$
($4\,Dq$) en daarna voor $d^3$ en $d^8$ ($2\,Dq$), nul voor $d^0$, hoogspin-$d^5$ en $d^{10}$, en negatief
voor hoogspin-$d^4$ en $d^9$.
\end{proposition}

\begin{proof}
Laat elk ligand een $d$-orbitaal verhogen met $e_\sigma$ maal het kwadraat van de
hoekamplitude van de orbitaal in zijn richting, genormeerd op 1 langs de
as van de lob: een axiaal
ligand geeft $d_{z^2}$ 1, een equatoriaal ligand geeft $d_{z^2}$ $\frac14$ en $d_{x^2-y^2}$ $\frac34$; de $t_{2g}$-orbitalen wijzen
tussen de liganden door en krijgen niets. De octaëder geeft $d_{z^2}$ en $d_{x^2-y^2}$ elk $3e_\sigma$
($\Delta_{\mathrm o} = 3e_\sigma = 10Dq$); de vierkante piramide, zonder één axiaal ligand, geeft $d_{z^2}$ $2e_\sigma$ en
$d_{x^2-y^2}$ $3e_\sigma$. De stabilisatie van $n$ elektronen is hun energie min $n/5$ maal de som van
de vijf niveaus (het zwaartepunt). Voor laagspin-$d^6$ ($t_{2g}^6$): $0 - \frac65 \times 6e_\sigma$ in de
octaëder, $0 - \frac65 \times 5e_\sigma$ in de piramide: een verlies van $1.2e_\sigma = 4Dq$. Voor $d^3$: $\frac35e_\sigma =
2Dq$; voor $d^8$ ($t_{2g}^6e_g^2$): $(5 - \frac85 \times 5) - (6 - \frac85 \times 6) = 0.6e_\sigma$. Voor $d^9$: $(7 - 9) - (9 -
10.8) = -0.2e_\sigma$.
\end{proof}

\begin{omfigure}
\begin{tikzpicture}
\begin{axis}[width=0.72\linewidth, height=5.4cm, xmin=-0.6, xmax=10.6, ymin=-1.2, ymax=4.6,
  axis x line=bottom, axis y line=left, ycomb, xtick={0,1,2,3,4,5,6,7,8,9,10},
  xlabel={$n$ ($d^n$)}, ylabel={activeringsenergie / $Dq$}, tick label style={font=\small},
  label style={font=\small}]
\draw[black!40] (axis cs:-0.6,0) -- (axis cs:10.6,0);
\addplot[omDef, line width=4pt, xshift=-2.5pt] table[x=n, y=hs] {figdata/out/bachelor-3-ligand-field-activation.dat};
\addplot[omThm, line width=4pt, xshift=2.5pt, unbounded coords=discard] table[x=n, y=ls] {figdata/out/bachelor-3-ligand-field-activation.dat};
\end{axis}
\end{tikzpicture}

{\small Ligandveldactiveringsenergieën voor een dissociatieve weg via een
vierkante piramide, in het model met alleen $\sigma$-binding van de propositie (blauw: hoogspin;
rood: laagspin). De grootste waarden,
voor laagspin-$d^6$, $d^3$ en $d^8$, passen bij de inerte of trage ionen (\ce{Co^{3+}}, \ce{Cr^{3+}}, \ce{Ni^{2+}});
de negatieve, voor $d^4$ en $d^9$, bij de zeer labiele jahn-tellerionen.}
\end{omfigure}

\begin{definition}[Activeringsvolume]\label{def:b3:complex-mechanisms:activation-volume}
Het \emph{activeringsvolume}\index{activeringsvolume} $\Delta V^{\ddagger}$ van een elementaire
stap is het verschil tussen het partieel molair volume van het
\omterm{def:b3:rate-theories:activated-complex}{geactiveerde complex} en die van de beginstoffen.
\end{definition}

\begin{proposition}[Drukafhankelijkheid van een snelheidsconstante]\label{prop:b3:complex-mechanisms:pressure}
Bij constante temperatuur is $\displaystyle\Bigl(\frac{\partial\ln k}{\partial p}\Bigr)_T = -\frac{\Delta V^{\ddagger}}{RT}$.
\end{proposition}

\begin{proof}
Volgens de \omterm{thm:b3:rate-theories:eyring}{vergelijking van Eyring} is $\ln k = \text{constante} - \Delta^{\ddagger}G/RT$, en $(\partial\Delta^{\ddagger}G/\partial p)_T = \Delta V^{\ddagger}$, zoals voor
elke gibbsenergie.
\end{proof}

Een stap waarin een ligand vertrekt, heeft $\Delta V^{\ddagger} > 0$; een stap waarin een ligand binnenkomt, heeft $\Delta V^{\ddagger}
< 0$. Gemeten bij enkele honderden megapascal is het \omterm{def:b3:complex-mechanisms:activation-volume}{activeringsvolume} de
meest directe diagnose van een substitutiemechanisme, minder vertroebeld
door solvatatie dan de \omterm{def:b3:rate-theories:activation}{activeringsentropie}.

\section{Substitutiemechanismen}

\begin{definition}[Dissociatief, associatief en uitwisselingsmechanisme]\label{def:b3:complex-mechanisms:mechanisms}
Bij een \emph{dissociatief mechanisme}\index{dissociatief mechanisme} (D)
vertrekt het uittredende ligand eerst, wat een intermediair met een
lager coördinatiegetal geeft; bij een
\emph{associatief mechanisme}\index{associatief mechanisme} (A) bindt het
binnenkomende ligand eerst, wat een intermediair met een hoger
coördinatiegetal geeft; bij een
\emph{uitwisselingsmechanisme}\index{uitwisselingsmechanisme} (I) wisselen
de twee liganden in één stap, met een karakter van bindingsbreuk ($I_d$) of bindingsvorming ($I_a$).
\end{definition}

\begin{proposition}[Snelheidswetten van de mechanismen D en A]\label{prop:b3:complex-mechanisms:d-a-laws}
Voor $\mathrm{ML_5X + Y \to ML_5Y + X}$ geeft een mechanisme D ($\mathrm{ML_5X} \rightleftharpoons \mathrm{ML_5} + \mathrm X$, $k_1$, $k_{-1}$; $\mathrm{ML_5} + \mathrm Y
\to \mathrm{ML_5Y}$, $k_2$)
\[
  v = \frac{k_1k_2[\mathrm{ML_5X}][\mathrm Y]}{k_{-1}[\mathrm X] + k_2[\mathrm Y]},
\]
van de eerste orde in het complex en onafhankelijk van $[\mathrm Y]$ bij hoge $[\mathrm Y]$; een mechanisme A
via een zevengecoördineerd intermediair in een stationaire toestand geeft $v =
k[\mathrm{ML_5X}][\mathrm Y]$, van de eerste orde in elk.
\end{proposition}

\begin{proof}
D: stationaire toestand voor $\mathrm{ML_5}$, $k_1[\mathrm{ML_5X}] = (k_{-1}[\mathrm X] + k_2[\mathrm Y])[\mathrm{ML_5}]$, en $v = k_2[\mathrm{ML_5}][\mathrm Y]$.
Bij hoge $[\mathrm Y]$ gaat $v \to k_1[\mathrm{ML_5X}]$. A: $\mathrm{ML_5X} + \mathrm Y \rightleftharpoons \mathrm{ML_5XY}$ ($k_a$, $k_{-a}$), $\mathrm{ML_5XY} \to
\mathrm{ML_5Y} + \mathrm X$ ($k_b$); de stationaire toestand geeft $v = \frac{k_ak_b}{k_{-a} + k_b}[\mathrm{ML_5X}][\mathrm Y]$.
\end{proof}

In water is het binnenkomende ligand meestal in veel lagere concentratie
aanwezig dan het oplosmiddel, en de eigenlijke eerste stap is de
ontmoeting van het complex en het ligand.

\begin{definition}[Mechanisme van Eigen-Wilkins]\label{def:b3:complex-mechanisms:eigen-wilkins}
Het \emph{mechanisme van Eigen-Wilkins}\index{mechanisme van Eigen-Wilkins} van
de substitutie van een aquacomplex is een snel voorevenwicht dat een
\emph{buitensfeercomplex}\index{buitensfeercomplex} vormt, waarin het
binnenkomende ligand naast het complex zit, buiten zijn
coördinatiesfeer, gevolgd door de snelheidsbepalende uitwisseling van een
watermolecuul tegen dat ligand.
\end{definition}

\begin{theorem}[Snelheidswet van Eigen-Wilkins]\label{thm:b3:complex-mechanisms:eigen-wilkins}
Met de evenwichtsconstante van de buitensfeer $K_{\mathrm{os}}$ en de snelheidsconstante van de uitwisseling
$k_i$ is de waargenomen pseudo-eersteordesnelheidsconstante bij overmaat $[\mathrm Y]$
\[
  k_{\mathrm{obs}} = \frac{K_{\mathrm{os}}k_i[\mathrm Y]}{1 + K_{\mathrm{os}}[\mathrm Y]} .
\]
\end{theorem}

\begin{proof}
Het complex is verdeeld over de vrije vorm en de buitensfeervorm, $[\mathrm{M\cdot Y}] = K_{\mathrm{os}}[\mathrm M][\mathrm Y]$,
met totaal $[\mathrm M]_{\mathrm{tot}} = [\mathrm M](1 + K_{\mathrm{os}}[\mathrm Y])$; de snelheid is $k_i[\mathrm{M\cdot Y}] = k_iK_{\mathrm{os}}[\mathrm Y][\mathrm M]_{\mathrm{tot}}/(1 + K_{\mathrm{os}}[\mathrm Y])$.
\end{proof}

Bij lage $[\mathrm Y]$ is de reactie van de tweede orde met $k = K_{\mathrm{os}}k_i$; $K_{\mathrm{os}}$, geschat uit de
ladingen en de kleinste naderingsafstand, is vrijwel dezelfde voor alle
liganden met dezelfde lading, zodat $k_i$, de stap van de wateruitwisseling, de
snelheid bepaalt: substitutie op een aqua-ion is vrijwel onafhankelijk
van het binnenkomende ligand, de handtekening van een dissociatieve
uitwisseling.

\begin{definition}[Mechanisme van de geconjugeerde base]\label{def:b3:complex-mechanisms:conjugate-base}
Het \emph{mechanisme van de geconjugeerde base}\index{mechanisme van de geconjugeerde base}
van de basische hydrolyse van een amminecomplex, $\mathrm{[Co(NH_3)_5X]^{2+} + OH^- \to [Co(NH_3)_5(OH)]^{2+} + X^-}$, is
de deprotonering van een ammineligand in een snel voorevenwicht, gevolgd
door het dissociatieve verlies van $\mathrm X^-$ uit de geconjugeerde amidobase, waarvan het ligand $\mathrm{NH_2^-}$
het complex \omterm{def:b3:complex-mechanisms:labile}{labiel} maakt.
\end{definition}

\begin{proposition}[Snelheidswet van de basische hydrolyse]\label{prop:b3:complex-mechanisms:conjugate-base}
Met de deprotoneringsconstante $K$ en de dissociatiesnelheidsconstante $k$
van de geconjugeerde base is $v = \dfrac{kK[\text{complex}]_{\mathrm{tot}}[\ce{OH-}]}{1 + K[\ce{OH-}]}$, van de tweede orde bij lage $[\ce{OH-}]$.
\end{proposition}

\begin{proof}
Zoals voor de wet van Eigen-Wilkins: het complex is verdeeld over zijn
zure vorm en zijn geconjugeerde base in de verhouding $1 : K[\ce{OH-}]$, en alleen de laatste reageert.
\end{proof}

De snelheidswet alleen onderscheidt dit mechanisme niet van een
rechtstreekse aanval van
\ce{OH-}; het bewijs is dat basische hydrolyse een N--H-proton vereist (complexen zonder
zo'n proton vertonen haar niet) en dat de protonen van de ammine veel
sneller met het oplosmiddel uitwisselen dan de hydrolyse verloopt.

\begin{method}[Een substitutiemechanisme diagnosticeren]\label{met:b3:complex-mechanisms:diagnose}
\begin{enumerate}
  \item Afhankelijkheid van de binnenkomende groep: snelheid onafhankelijk
        van Y (of verzadigend): dissociatief; snelheid evenredig met $[\mathrm Y]$ en gevoelig voor de aard van Y:
        associatief.
  \item \omterm{def:b3:complex-mechanisms:activation-volume}{Activeringsvolume}: positief, dissociatief; negatief, associatief.
  \item \omterm{def:b3:rate-theories:activation}{Activeringsentropie}: positief voor D, negatief voor A (met
        voorzichtigheid, want de solvatatie draagt bij).
  \item Afhankelijkheid van de uittredende groep: een snelheid die de
        sterkte van de binding M--X volgt, wijst op bindingsbreuk in de
        overgangstoestand.
\end{enumerate}
\end{method}

\begin{omfigure}
\begin{tikzpicture}
\begin{axis}[name=pd, omprofile, width=0.46\linewidth, height=4.6cm, xmin=0, xmax=10, ymin=0, ymax=10,
  xlabel={reactiecoördinaat}, ylabel={energie}, title={\small D: $\mathrm{ML_5X} \to \mathrm{ML_5} + \mathrm X \to \mathrm{ML_5Y}$}]
\addplot[omDef, very thick, smooth] coordinates {(0,2) (1.5,2.1) (3,7.5) (4.2,5.6) (5.0,5.4) (5.8,5.6) (7,7.0) (8.5,1.6) (10,1.5)};
\node[font=\scriptsize, anchor=north, align=center] at (axis cs:5,5.2) {ML$_5$ (CG 5)};
\end{axis}
\begin{axis}[at={(pd.east)}, anchor=west, xshift=1.0cm, omprofile, width=0.46\linewidth, height=4.6cm,
  xmin=0, xmax=10, ymin=0, ymax=10, xlabel={reactiecoördinaat}, ylabel={energie},
  title={\small A: $\mathrm{ML_5X} + \mathrm Y \to \mathrm{ML_5XY} \to \mathrm{ML_5Y}$}]
\addplot[omThm, very thick, smooth] coordinates {(0,2) (1.5,2.1) (3,6.5) (4.2,4.4) (5.0,4.2) (5.8,4.4) (7,7.2) (8.5,1.6) (10,1.5)};
\node[font=\scriptsize, anchor=north, align=center] at (axis cs:5,4.0) {ML$_5$XY (CG 7)};
\end{axis}
\end{tikzpicture}

{\small Energieprofielen van een dissociatieve en een associatieve
substitutie (schematisch): beide lopen via een intermediair, met een lager
coördinatiegetal (CG) bij D en een hoger coördinatiegetal bij A. In de
\omterm{def:b3:complex-mechanisms:mechanisms}{uitwisselingsmechanismen} verdwijnt de put en blijft één enkele
overgangstoestand over.}
\end{omfigure}

Een octaëdrisch complex dat via een vierkante piramide substitueert,
behoudt in het algemeen zijn configuratie ($cis$ blijft $cis$); een complex dat via een trigonale
bipiramide verloopt, zoals veel kobalt(III)complexen bij de basische
hydrolyse, kan een
mengsel van producten $cis$ en $trans$ geven.

\section{Substitutie in vlakke vierkante complexen en het trans-effect}

Vlakke vierkante complexen van $d^8$-ionen (\ce{Pt^{2+}}, \ce{Pd^{2+}}, \ce{Au^{3+}}, \ce{Ni^{2+}} met sterke liganden)
zijn coördinatief onverzadigd: een binnenkomend ligand kan naderen langs
de lege axiale richting. Hun substituties zijn associatief, via een
vijfgecoördineerde trigonaal-bipiramidale overgangstoestand.

\begin{proposition}[Snelheidswet met twee termen]\label{prop:b3:complex-mechanisms:square-planar}
De substitutie $\mathrm{[PtL_3X] + Y \to [PtL_3Y] + X}$ in een coördinerend oplosmiddel S volgt
$k_{\mathrm{obs}} = k_1 + k_2[\mathrm Y]$ bij overmaat Y.
\end{proposition}

\begin{proof}
Twee parallelle associatieve wegen: rechtstreekse aanval van Y ($k_2[\mathrm Y]$), en aanval van het
oplosmiddel, dat in grote overmaat aanwezig is (pseudo-eerste orde $k_1$), wat $[\mathrm{PtL_3S}]$ geeft, dat
daarna snel door Y omgezet wordt. Parallelle wegen van de eerste orde
tellen op.
\end{proof}

\begin{definition}[Trans-effect, trans-invloed]\label{def:b3:complex-mechanisms:trans-effect}
Het \emph{trans-effect}\index{trans-effect} is de versnelling van de
substitutie van een ligand door het ligand dat er trans tegenover staat:
een kinetisch effect, op de overgangstoestand.
De \emph{trans-invloed}\index{trans-invloed} is de verzwakking, in de
grondtoestand, van de binding trans ten opzichte van een ligand, zichtbaar
in bindingslengten en vibratiefrequenties: een thermodynamisch effect.
\end{definition}

Het \omterm{def:b3:complex-mechanisms:trans-effect}{trans-effect} volgt de volgorde $\ce{CO}$, $\ce{CN-}$, $\ce{C2H4} > \ce{PR3}$, $\ce{H-} > \ce{CH3-} > \ce{I-}$, $\ce{SCN-} >
\ce{Br-} > \ce{Cl-} > \ce{NH3}$, $\ce{H2O}$. Sterke $\sigma$-donoren verzwakken de transbinding (\omterm{def:b3:complex-mechanisms:trans-effect}{trans-invloed});
sterke $\pi$-acceptoren stabiliseren de vijfgecoördineerde overgangstoestand door
\omterm{def:b3:computational-chemistry:dft}{elektronendichtheid} aan het metaal te onttrekken. Beide verhogen de
snelheid.

\begin{method}[De synthese van isomeren van platina(II) plannen]\label{met:b3:complex-mechanisms:pt-synthesis}
\begin{enumerate}
  \item Bij elke stap wordt het ligand vervangen dat trans staat ten
        opzichte van het aanwezige ligand met het grootste \omterm{def:b3:complex-mechanisms:trans-effect}{trans-effect}.
  \item Orden de toevoegingen zo dat deze regel naar het gewenste isomeer
        leidt.
\end{enumerate}
\end{method}

\begin{omfigure}
\setchemfig{atom sep=2.2em}
\schemestart
\chemfig{Cl-Pt(-[2]Cl)(-[6]Cl)-Cl}
\arrow{->[\ce{NH3}]}
\chemfig{Cl-Pt(-[2]Cl)(-[6]Cl)-NH_3}
\arrow{->[\ce{NH3}]}
\chemfig{Cl-Pt(-[2]NH_3)(-[6]Cl)-NH_3}
\schemestop

\bigskip
\schemestart
\chemfig{H_3N-Pt(-[2]NH_3)(-[6]NH_3)-NH_3}
\arrow{->[\ce{Cl-}]}
\chemfig{H_3N-Pt(-[2]NH_3)(-[6]NH_3)-Cl}
\arrow{->[\ce{Cl-}]}
\chemfig{Cl-Pt(-[2]NH_3)(-[6]NH_3)-Cl}
\schemestop

{\small Het \omterm{def:b3:complex-mechanisms:trans-effect}{trans-effect} aan het werk (ladingen weggelaten). Boven: vanuit $\ce{[PtCl4]^{2-}}$ vervangt de tweede
ammoniak een chloride trans ten opzichte van chloride (Cl$^-$ heeft het grotere \omterm{def:b3:complex-mechanisms:trans-effect}{trans-effect}),
wat het isomeer $cis$ geeft, cisplatine. Onder: vanuit $\ce{[Pt(NH3)4]^{2+}}$ vervangt het tweede chloride
de ammoniak trans ten opzichte van het eerste chloride, wat het isomeer $trans$ geeft.}
\end{omfigure}

\section{Elektronenoverdracht}

\begin{definition}[Elektronenoverdracht via de buitensfeer en via de binnensfeer]\label{def:b3:complex-mechanisms:electron-transfer}
Bij
\emph{elektronenoverdracht via de buitensfeer}\index{elektronenoverdracht via de buitensfeer}
gaat het elektron over tussen twee complexen waarvan de
coördinatiesferen intact blijven; bij
\emph{elektronenoverdracht via de binnensfeer}\index{elektronenoverdracht via de binnensfeer}
gaat het via een ligand dat de twee metalen overbrugt in een kortlevend
tweekernig complex. Een
\emph{zelfuitwisselingsreactie}\index{zelfuitwisselingsreactie} is een
elektronenoverdracht tussen de twee oxidatietoestanden van hetzelfde
koppel, zonder netto chemische verandering, zoals $\ce{[Fe(H2O)6]^{3+} + [Fe(H2O)6]^{2+}}$.
\end{definition}

Henry Taube toonde de weg via de binnensfeer in 1953 aan met een reactie
die zo gekozen was dat het product zijn verhaal vertelde. Kobalt(III) is
inert, kobalt(II) \omterm{def:b3:complex-mechanisms:labile}{labiel}; chroom(II)
is \omterm{def:b3:complex-mechanisms:labile}{labiel}, chroom(III) inert. Wanneer $\ce{[Co(NH3)5Cl]^{2+}}$ gereduceerd wordt door $\ce{[Cr(H2O)6]^{2+}}$ in
zuur milieu, draagt al het gevormde chroom(III) het chloride, als $\ce{[Cr(H2O)5Cl]^{2+}}$, zelfs in
aanwezigheid van vrij chloride: het chloride moet de twee metalen
overbrugd hebben, gebonden zijn aan het inerte chroom(III) zodra het
elektron overging, en losgelaten zijn door het labiele kobalt(II).

\begin{omfigure}
\begin{tikzpicture}[font=\small]
  \begin{scope}[scale=0.8]
    \pic[transform shape] (a) at (0,0) {omoctahedron};
    \pic[transform shape] (b) at (2.0,0) {omoctahedron};
  \end{scope}
  \node[anchor=north east] at (0.3,-1.0) {$\mathrm{Co^{III}(NH_3)_5}$};
  \node[anchor=north west] at (1.3,-1.0) {$\mathrm{Cr^{II}(H_2O)_5}$};
  \node[omThm, anchor=south] at (0.8,0.15) {Cl};
  \draw[->, thick, omDef] (1.6,1.75) to[bend right=40] node[midway, above] {e$^-$} (0,1.75);
\end{tikzpicture}

{\small Het overbrugde intermediair van Taube (schematisch): het chloride,
nog gebonden aan kobalt(III), treedt de coördinatiesfeer van chroom(II)
binnen door een watermolecuul te vervangen; het elektron gaat over via de
brug, en het chloride blijft op het nu inerte chroom(III).}
\end{omfigure}

Een elektron springt in ongeveer $10^{-16}$ s, veel sneller dan kernen bewegen: volgens
het \omterm{thm:b3:electronic-spectroscopy:franck-condon}{franck-condonprincipe} (\cref{ch:b3:electronic-spectroscopy}) kan de
overdracht alleen plaatsvinden bij een kernconfiguratie waarin
beginstoffen en producten dezelfde energie hebben. Bij de
zelfuitwisseling van ijzer zijn de bindingen \ce{Fe^{III}}--O korter dan
de bindingen \ce{Fe^{II}}--O; voordat het elektron kan bewegen, moeten de twee complexen vervormen tot
een gemeenschappelijke tussengeometrie, wat energie kost.

\section{Theorie van Marcus}

\begin{definition}[Reorganisatie-energie]\label{def:b3:complex-mechanisms:reorganisation}
De \emph{reorganisatie-energie}\index{reorganisatie-energie} $\lambda$ van
een elektronenoverdracht is de energie die nodig is om de beginstoffen,
zonder het elektron over te dragen, in de kernconfiguratie
(bindingslengten van de complexen en oriëntaties van het omringende
oplosmiddel) van de producten te brengen. Ze is de som van een inwendig
deel, van de bindingen, en een uitwendig deel, van het oplosmiddel.
\end{definition}

\begin{theorem}[Vergelijking van Marcus]\label{thm:b3:complex-mechanisms:marcus}
Als de gibbsenergieën van beginstoffen en producten parabolen met
dezelfde kromming zijn in een reactiecoördinaat $x$ (0 bij het evenwicht van de beginstoffen, 1 bij dat
van de producten), is de \omterm{def:b3:rate-theories:activation}{gibbsenergie van activering} van een
elektronenoverdracht met
standaardreactiegibbsenergie $\Delta G^\circ$ en \omterm{def:b3:complex-mechanisms:reorganisation}{reorganisatie-energie} $\lambda$
\[
  \Delta G^{\ddagger} = \frac{(\lambda + \Delta G^\circ)^2}{4\lambda} .
\]
Dit is de \emph{vergelijking van Marcus}\index{vergelijking van Marcus}.
\end{theorem}

\begin{proof}
Schrijf $G_R = \lambda x^2$ en $G_P = \lambda(x - 1)^2 + \Delta G^\circ$: de kromming ligt vast door $G_R(1) = \lambda$, de
definitie van $\lambda$. De overdracht vindt plaats waar de krommen elkaar snijden: $\lambda x^2 = \lambda x^2 - 2\lambda x + \lambda +
\Delta G^\circ$, dus $x^\ast = (\lambda + \Delta G^\circ)/2\lambda$, en $\Delta G^{\ddagger} = G_R(x^\ast) = (\lambda + \Delta G^\circ)^2/4\lambda$.
\end{proof}

\begin{corollary}[Geïnverteerd gebied]\label{cor:b3:complex-mechanisms:inverted}
Bij vaste $\lambda$ neemt de snelheid toe naarmate de reactie exergonischer wordt, tot
$-\Delta G^\circ = \lambda$, waar $\Delta G^{\ddagger} = 0$; daarna neemt ze weer af.
\end{corollary}

\begin{proof}
$\Delta G^{\ddagger}$ is een parabool in $\Delta G^\circ$ met haar minimum, nul, bij $\Delta G^\circ = -\lambda$; voor $\Delta G^\circ < -\lambda$ groeit ze
weer.
\end{proof}

\begin{definition}[Geïnverteerd gebied]\label{def:b3:complex-mechanisms:inverted}
Het \emph{geïnverteerd gebied}\index{geinverteerd gebied@geïnverteerd gebied}
van de elektronenoverdracht is het gebied van drijvende krachten $-\Delta G^\circ > \lambda$ waarin de snelheid afneemt naarmate de reactie
exergonischer wordt.
\end{definition}

\begin{omfigure}
\begin{tikzpicture}
\begin{axis}[name=mp, width=0.47\linewidth, height=5.6cm, xmin=-0.6, xmax=2.2, ymin=-160, ymax=120,
  axis lines=left, xlabel={reactiecoördinaat $x$}, ylabel={$G$ / \unit{kJ.mol^{-1}}},
  tick label style={font=\small}, label style={font=\small}, xtick={0,1,2}]
\addplot[black, thick] table[x=x, y=GR] {figdata/out/bachelor-3-marcus-parabolas.dat};
\addplot[omDef!50, thick] table[x=x, y=P0] {figdata/out/bachelor-3-marcus-parabolas.dat};
\addplot[omDef, thick] table[x=x, y=P50] {figdata/out/bachelor-3-marcus-parabolas.dat};
\addplot[omThm, thick] table[x=x, y=P100] {figdata/out/bachelor-3-marcus-parabolas.dat};
\addplot[omThm!60, thick, dashed] table[x=x, y=P150] {figdata/out/bachelor-3-marcus-parabolas.dat};
\node[font=\scriptsize, anchor=north] at (axis cs:0,-2) {R};
\end{axis}
\begin{axis}[at={(mp.east)}, anchor=west, xshift=1.4cm, width=0.47\linewidth, height=5.6cm,
  xmin=0, xmax=250, ymin=0, ymax=11.5, axis lines=left, xlabel={$-\Delta G^\circ$ / \unit{kJ.mol^{-1}}},
  ylabel={$\log(k/\unit{s^{-1}})$}, tick label style={font=\small}, label style={font=\small}]
\addplot[omDef, very thick] table[x=mdG, y=logk] {figdata/out/bachelor-3-marcus-rate.dat};
\draw[black!40, dashed] (axis cs:100,0) -- (axis cs:100,11);
\node[font=\scriptsize, anchor=south east] at (axis cs:95,1) {normaal};
\node[font=\scriptsize, anchor=south west] at (axis cs:105,1) {geïnverteerd};
\end{axis}
\end{tikzpicture}

{\small Theorie van Marcus met $\lambda = \qty{100}{kJ/mol}$ (model). Links: de parabool van de
beginstoffen (zwart) en de parabolen van de producten voor $\Delta G^\circ = 0$ en $-50$ (blauw), $-100 = -\lambda$ (rood) en
\qty{-150}{kJ/mol} (gestreept, \omterm{def:b3:complex-mechanisms:inverted}{geïnverteerd gebied}); hun snijpunt, de
overgangstoestand, zakt tot op de bodem van de kromme van de beginstoffen bij
$-\Delta G^\circ = \lambda$ en stijgt daarna weer. Rechts: de logaritme van de snelheidsconstante
(modelprefactor \qty{e11}{s^{-1}}) gaat door een maximum bij $-\Delta G^\circ = \lambda$.}
\end{omfigure}

\begin{proposition}[Reorganisatie-energie van de buitensfeer]\label{prop:b3:complex-mechanisms:outer-reorganisation}
Voor twee bolvormige beginstoffen met stralen $a_1$, $a_2$ op een afstand $d$ in een oplosmiddel met
brekingsindex $n$ en relatieve permittiviteit $\varepsilon_s$ is het oplosmiddeldeel van $\lambda$
evenredig met $\bigl(\frac{1}{2a_1} + \frac{1}{2a_2} - \frac1d\bigr)\bigl(\frac{1}{n^2} - \frac{1}{\varepsilon_s}\bigr)$.
\end{proposition}

\admitted

Water, zeer polair ($\varepsilon_s$ groot) maar met een gewone brekingsindex, geeft een grote
reorganisatie van het oplosmiddel; apolaire oplosmiddelen een kleine.
Grote beginstoffen reorganiseren minder: daarom begraven
elektronenoverdrachtseiwitten hun metaalcentra binnen in het eiwit.

\begin{theorem}[Kruisrelatie van Marcus]\label{thm:b3:complex-mechanisms:cross-relation}
Voor een kruisreactie $\mathrm{A_{ox} + B_{red} \to A_{red} + B_{ox}}$ met evenwichtsconstante $K_{12}$, tussen koppels
waarvan de zelfuitwisselingssnelheidsconstanten $k_{11}$ en $k_{22}$ zijn, geldt
\[
  k_{12} = \sqrt{k_{11}k_{22}K_{12}f}, \qquad \ln f = \frac{(\ln K_{12})^2}{4\ln(k_{11}k_{22}/Z^2)},
\]
met $Z$ de botsingsfrequentiefactor; $f \approx 1$ wanneer $K_{12}$ niet te groot is. Dit is de
\emph{kruisrelatie van Marcus}\index{kruisrelatie van Marcus}.
\end{theorem}

\begin{proof}[Gedeeltelijk bewijs]
Neem aan dat de \omterm{def:b3:complex-mechanisms:reorganisation}{reorganisatie-energie} van de kruisreactie het gemiddelde
is van die van de
zelfuitwisselingen, $\lambda_{12} = (\lambda_{11} + \lambda_{22})/2$ (elke beginstof draagt de helft van zijn eigen
reorganisatie bij de zelfuitwisseling bij), en alle snelheidsconstanten $k = Z\eu^{-\Delta G^{\ddagger}/RT}$. Dan is
$\Delta G_{12}^{\ddagger} = \frac{\lambda_{12}}{4}\bigl(1 + \frac{\Delta G^\circ_{12}}{\lambda_{12}}\bigr)^2 = \frac{\lambda_{12}}{4} + \frac{\Delta G^\circ_{12}}{2} + \frac{(\Delta G^\circ_{12})^2}{4\lambda_{12}}$. De eerste term geeft
$\sqrt{k_{11}k_{22}}$ (want $\Delta G^{\ddagger}_{ii} = \lambda_{ii}/4$), de tweede $\sqrt{K_{12}}$ (met $\Delta G^\circ_{12} = -RT\ln K_{12}$), de derde
de factor $f$.
\end{proof}

\begin{method}[De snelheid van een kruisreactie uit zelfuitwisselingsgegevens]\label{met:b3:complex-mechanisms:marcus-estimate}
\begin{enumerate}
  \item Zoek de zelfuitwisselingsconstanten $k_{11}$, $k_{22}$ van de twee koppels.
  \item Bereken $K_{12}$ uit de standaardpotentialen: $\ln K_{12} = nF(E^\circ_{\text{oxidator}} - E^\circ_{\text{reductor}})/RT$.
  \item $k_{12} = \sqrt{k_{11}k_{22}K_{12}}$ (met $f$ als $K_{12}$ zeer groot is). Overeenstemming binnen een
        kleine factor met de gemeten waarde wijst op een mechanisme via
        de buitensfeer.
\end{enumerate}
\end{method}

Marcus publiceerde de theorie in 1956; het geïnverteerde gebied, eerst
betwijfeld, werd in 1984 waargenomen door Gerhard Closs en John Miller in
moleculen waarin donor en acceptor door een starre afstandhouder op een
vaste afstand gehouden worden, zodat diffusie het niet kan verbergen.

\begin{inthelab}[Wateruitwisseling met zuurstof-17-NMR]
De zuurstof-17-resonantie van water dat aan een paramagnetisch ion
gebonden is en die van het vrije water worden verbreed door de
uitwisseling ertussen. Het registreren van de lijnbreedte van het signaal
van het vrije water als functie van de temperatuur, in een oplossing van
het metaalperchloraat
verrijkt met \ce{^{17}O}, geeft de uitwisselingssnelheidsconstante en, in een
hogedrukprobe, haar \omterm{def:b3:complex-mechanisms:activation-volume}{activeringsvolume}.
\end{inthelab}

\begin{safety}{\ghs[0.8cm]{GHS05}\ghs[0.8cm]{GHS06}\ghs[0.8cm]{GHS08}}
% ledger: ghs4:K2PtCl4
Kaliumtetrachloroplatinaat(II): giftig bij inslikken, veroorzaakt ernstig
oogletsel, sensibiliserend voor de luchtwegen en de huid.
Platinacomplexen worden met handschoenen in een zuurkast gehanteerd;
cisplatine en zijn analogen zijn cytotoxische geneesmiddelen en worden
als zodanig behandeld.
\end{safety}

\begin{history}[Taube en Marcus]
Henry Taube deelde complexen in 1952 in als \omterm{def:b3:complex-mechanisms:labile}{labiel} of inert, en stelde
met zijn experiment met chroom en kobalt van 1953 het mechanisme via de
binnensfeer vast; hij ontving de Nobelprijs voor de Scheikunde van 1983.
Rudolph Marcus leidde vanaf 1956 de theorie van de elektronenoverdracht
af die zijn naam draagt, en ontving de Nobelprijs voor de Scheikunde van
1992.
\end{history}

\section{Oefeningen}

\begin{exercise}[$\star$]\label{exo:b3:complex-mechanisms:1}
Deel in als \omterm{def:b3:complex-mechanisms:labile}{labiel} of inert, met een reden: $\ce{[Cr(H2O)6]^{3+}}$, $\ce{[Co(NH3)6]^{3+}}$, $\ce{[Cr(H2O)6]^{2+}}$,
$\ce{[Cu(H2O)6]^{2+}}$, $\ce{[Zn(H2O)6]^{2+}}$.
\end{exercise}

\begin{exercise}[$\star$]\label{exo:b3:complex-mechanisms:2}
Toon aan dat de snelheidswet van D herleid wordt tot $v = k_1[\mathrm{ML_5X}]$ bij hoge $[\mathrm Y]$ en tot $v = (k_1k_2/k_{-1})
[\mathrm{ML_5X}][\mathrm Y]/[\mathrm X]$ als $k_{-1}[\mathrm X] \gg k_2[\mathrm Y]$.
\end{exercise}

\begin{exercise}[$\star$]\label{exo:b3:complex-mechanisms:3}
Voorspel de tekens van $\Delta V^{\ddagger}$ en $\Delta^{\ddagger}S^\circ$ voor een substitutie D en een substitutie A.
\end{exercise}

\begin{exercise}[$\star$]\label{exo:b3:complex-mechanisms:4}
Geef de producten van $\ce{[PtCl4]^{2-}}$ met twee \ce{NH3}, en van $\ce{[Pt(NH3)4]^{2+}}$ met twee \ce{Cl-}.
\end{exercise}

\begin{exercise}[$\star\star$]\label{exo:b3:complex-mechanisms:5}
Voor een substitutie op een aqua-ion is $K_{\mathrm{os}} = \qty{2.0}{L.mol^{-1}}$ en $k_i = \qty{3.0e4}{s^{-1}}$ (gegevens van
de oefening). Bereken $k_{\mathrm{obs}}$ bij $[\mathrm Y] = 0.10$ en \qty{1.0}{mol/L}, en de limiet ervan.
\end{exercise}

\begin{exercise}[$\star\star$]\label{exo:b3:complex-mechanisms:6}
Een snelheidsconstante verdubbelt wanneer de druk stijgt van 0.1 tot \qty{100}{MPa} bij \qty{298}{K}.
Bereken $\Delta V^{\ddagger}$ en trek een besluit.
\end{exercise}

\begin{exercise}[$\star\star$]\label{exo:b3:complex-mechanisms:7}
In $\ce{[PtCl3(PEt3)]-}$ is de binding Pt--Cl trans ten opzichte van \ce{PEt3} \qty{2.42}{\angstrom} lang en zijn de twee andere \qty{2.32}{\angstrom}
(gegevens van de oefening). Interpreteer dit, en voorspel welk chloride
als eerste vervangen wordt.
\end{exercise}

\begin{exercise}[$\star\star$]\label{exo:b3:complex-mechanisms:8}
Som de experimentele aanwijzingen op die zouden aantonen dat een
\omterm{def:b3:complex-mechanisms:electron-transfer}{elektronenoverdracht via de binnensfeer} verloopt.
\end{exercise}

\begin{exercise}[$\star\star$]\label{exo:b3:complex-mechanisms:9}
Neem $\lambda = \qty{80}{kJ/mol}$ en bereken $\Delta G^{\ddagger}$ voor $\Delta G^\circ = 0$, $-20$, $-80$ en \qty{-120}{kJ/mol}.
\end{exercise}

\begin{exercise}[$\star\star\star$]\label{exo:b3:complex-mechanisms:10}
Twee koppels hebben de zelfuitwisselingsconstanten $k_{11} = 4.0$ en $k_{22} = \qty{2.0e5}{L.mol^{-1}.s^{-1}}$, en
hun kruisreactie heeft $K_{12} = \num{1.0e4}$ (gegevens van de oefening). Schat $k_{12}$ ($f = 1$).
\end{exercise}

\begin{exercise}[$\star\star\star$]\label{exo:b3:complex-mechanisms:11}
Bereken met het model met alleen $\sigma$-binding de
ligandveldactiveringsenergieën van ionen hoogspin-$d^5$, $d^7$ en $d^8$ en rangschik \ce{Mn^{2+}}, \ce{Co^{2+}} en \ce{Ni^{2+}} naar
de verwachte snelheid van wateruitwisseling.
\end{exercise}

\begin{exercise}[$\star\star\star$]\label{exo:b3:complex-mechanisms:12}
Waarom is het geïnverteerde gebied moeilijk waar te nemen bij reacties
tussen ionen die in oplossing naar elkaar toe diffunderen, en hoe hebben
starre donor-acceptormoleculen het zichtbaar gemaakt?
\end{exercise}

\section{Probleem: cisplatine maken, geen transplatine}

\begin{problem}[{Weekendprobleem --- cisplatine maken, geen transplatine:
het vlakke vierkante platina(II)ion, een synthese gepland met het
trans-effect, de aquatisering van het geneesmiddel, en waarom het binnen
de cellen geactiveerd wordt}]
\label{pb:b3:complex-mechanisms:1}
% ledger: ghs4:K2PtCl4
Cisplatine, $cis$-$\ce{[PtCl2(NH3)2]}$, wordt gemaakt uit $\ce{K2[PtCl4]}$. Gegevens van het probleem: bij
\qty{37}{\celsius} heeft de eerste aquatisering, $\ce{[PtCl2(NH3)2] + H2O -> [PtCl(H2O)(NH3)2]+ + Cl-}$, de
snelheidsconstante $k = \qty{9.0e-5}{s^{-1}}$ en de evenwichtsconstante $K = \qty{3.0e-3}{mol/L}$; de
chlorideconcentratie is ongeveer \qty{100}{mM} in bloedplasma en \qty{4}{mM} binnen cellen.

\medskip\noindent\textbf{Deel I --- Platina(II).}
\begin{enumerate}
  \item Geef het aantal $d$-elektronen van \ce{Pt^{2+}}.
  \item Waarom zijn zijn complexen vlak vierkant en diamagnetisch?
  \item Waarom zijn hun substituties associatief?
  \item Schrijf de snelheidswet met twee termen op en verklaar de twee
        termen.
  \item Is het geneesmiddel \omterm{def:b3:complex-mechanisms:labile}{labiel} of inert op de tijdschaal van een dag?
\end{enumerate}

\medskip\noindent\textbf{Deel II --- De synthese.}
\begin{enumerate}[resume]
  \item Wat geeft $\ce{K2[PtCl4]}$ met twee \ce{NH3}? Waarom?
  \item Die rechtstreekse route geeft een onzuiver product. De route van
        Dhara zet eerst
        $\ce{[PtCl4]^{2-}}$ om in $\ce{[PtI4]^{2-}}$ met overmaat KI. Waarom is jodide nuttig?
  \item Toevoegen van \ce{NH3} geeft $cis$-$\ce{[PtI2(NH3)2]}$. Verklaar de stereochemie.
  \item De jodiden worden daarna verwijderd met zilvernitraat in water.
        Wat ontstaat er?
  \item Toevoegen van KCl geeft cisplatine. Waarom verandert de
        configuratie niet?
  \item Hoe zou je in plaats daarvan transplatine maken?
  \item Waarom is transplatine onbruikbaar als geneesmiddel, hoewel het
        ook reageert?
\end{enumerate}

\medskip\noindent\textbf{Deel III --- Aquatisering.}
\begin{enumerate}[resume]
  \item Bereken de halveringstijd van de eerste aquatisering.
  \item Bereken de fractie die na \qty{2.0}{h} geaquatiseerd is in een
        chloridevrije oplossing.
  \item Waarom is het aquacomplex de reactieve vorm tegenover DNA?
  \item Is de aquatisering dissociatief of associatief? Wat zou haar
        \omterm{def:b3:complex-mechanisms:activation-volume}{activeringsvolume} zijn?
  \item Waarom wordt een oplossing van het geneesmiddel voor injectie in
        fysiologische zoutoplossing bereid?
\end{enumerate}

\medskip\noindent\textbf{Deel IV --- Plasma en cel.}
\begin{enumerate}[resume]
  \item Schrijf de evenwichtsfractie van het aquacomplex als functie van
        $[\ce{Cl-}]$.
  \item Bereken ze in plasma.
  \item Bereken ze binnen een cel.
  \item Bereken de verhouding.
  \item Waarom werkt het geneesmiddel daarom vooral binnen de cellen?
  \item Welke andere liganden in een cel kunnen platina binden en het
        geneesmiddel deactiveren?
  \item Waarom blijven de twee ammineliganden de hele tijd gebonden?
  \item Formuleer het resultaat: de verhouding van de fractie aquacomplex
        binnen een cel tot die in plasma.
\end{enumerate}
\end{problem}
